\subsection{全$(1,3)$表示の矢と四次元モジュライ}
\label{subsec:h3-order2-full-moduli}
定理\ref{thm:h3-order2-complete-conic-admissibility}の八係数許容開集合を
$\mathcal H\subset\mathbb A^8_{a,b,c,d,e,f,g,h}$と書く。
従って$e\ne0$、$(a,d)\ne(0,0)$、$\Delta_{\mathrm{AS}}\ne0$であり、
対応するconic bundle総空間が滑らかである。

\begin{theorem}\label{thm:h3-order2-full-arrows}
二つの幾何学的点に対応する全intrinsic pointed同値は、共通次数尺度を除いて
次の生成元写像で、かつそのような写像だけで与えられる：
\[
 p\mapsto\alpha p',\quad u\mapsto\beta u',\quad
 w\mapsto\gamma w',\quad v\mapsto\delta v'+r u'+s w',
 \qquad\alpha\beta\gamma\delta\ne0.
\]
係数の必要十分条件は
\begin{align*}
 \beta\delta a'&=a\alpha^2,&
 \beta\delta b'+2\beta r&=b\beta^2,&
 \beta\delta c'&=c\gamma^2,\\
 \gamma\delta d'&=d\alpha^2,&
 \gamma\delta e'&=e\beta^2,&
 \gamma\delta f'+2\gamma s&=f\gamma^2,\\
 ra'+sd'&=0,&&&\\
 \delta^2g'+\delta r b'+\delta s e'+r^2&=g\beta^2,&&&\\
 \delta^2h'+\delta r c'+\delta s f'+s^2&=h\gamma^2.&&&
\end{align*}
正規三次は$G\mapsto\alpha\gamma\delta G'$となる。
この係数群oidを共通次数尺度で割った表示stackは滑らか既約な
四次元Deligne--Mumford stackである。
さらに$ad\ne0$の開部分は、$a=d=e=1,b=0$とした
$\mathbb A^4_{c,f,g,h}$内の許容開集合と$B\mu_2$の直積である。
これは表示族のstackについての主張であり、任意のintrinsic Abel圏族の
認識・効果的降下や全H3内での成分性を含まない。
\end{theorem}
\begin{proof}
中山作用の正直線$kp$、混合関係から得る共通平面$\langle u,w\rangle$、
その上の二つの固有軸$ku,kw$は同型に自然である。
従って全幾何学的矢は表示した三角形に限られる。
六つの二次関係の係数を比較すると上の九式を得る。
特に最後の平方関係の$P'$係数は$\delta(ra'+sd')$である。
逆に九式は全関係を保ち、生成元行列は可逆なので全次数同型を与える。
正常三次の共変式は$\{u,w\}=0$と$[p,w]=0$で直接確認できる。
二次包絡の内在的復元により、商のpointed矢もこれらに尽くされる。

以下の相対stackは上の明示式を任意の基底環上で用いて定義する。
共通尺度を$\beta=1$で除く前に、sourceを固定するとtarget係数は
四つの可逆尺度と二つの平行移動から一意に決まり、残る条件は
\[
 ra/\beta+sd/\gamma=0
\]
だけである。$(a,d)$は単位イデアルを生成するのでこれは線束の全空間であり、
source写像は相対次元5の滑らかな写像である。
共通尺度を除くと相対次元4となる。合成は生成元写像の合成であり、
targetも同型により許容集合に属する。
非空開集合$\mathcal H$は滑らか既約八次元なので、商stackは滑らか既約四次元である。

安定化群のLie環を調べる。$\beta=1$、
$(\alpha,\gamma,\delta,r,s)=(1+A\epsilon,1+C\epsilon,1+D\epsilon,R\epsilon,S\epsilon)$
と置くと、$e\ne0$より
\[
 C=-D,\quad R=-bD/2,\quad S=-fD,
\]
および
\begin{gather*}
 a(D-2A)=dA=3cD=0,\
 (ab/2+df)D=(2g-b^2/2-ef)D=(4h-bc/2-f^2)D=0.
\end{gather*}
$ad\ne0$なら直ちに$A=D=0$。
$a=0,d\ne0$で$D\ne0$なら$c=f=h=0$となり$\Delta_{\mathrm{AS}}=0$。
$d=0,a\ne0$で$D\ne0$なら$b=c=0,g=ef/2,h=f^2/4$となり、やはり
$\Delta_{\mathrm{AS}}=0$である。従って全幾何学的安定化群はLie環零、
有限型アフィンなので有限で、標数0により有限エタールである。
この群oidの対角は非分岐、従ってstackはDMである。
対角の固有性・分離性はここでは主張しない。

$ad\ne0$ではsourceから正規形への矢を
\[
 \beta=1,\quad\alpha^2=e/d,\quad\gamma=d/a,\quad
 \delta=ae/d,\quad r=b/2,\quad s=-b/2
\]
で与え、九式から$c',f',g',h'$を一意に決められる。
必要な平方根はfppf局所的に存在する。
二つの正規形の間の矢は、九式から
$\gamma=\delta=1,r=s=0,\alpha^2=1$だけであり、残る四係数は不変である。
よって正規形開部分の群oidは恒等作用の$\mu_2$となり、主張を得る。
\end{proof}

\texttt{verify\_cubic\_order2\_full\_arrows.py}は普遍Clifford--Ore模型上で
全関係・三次共変式・正規化式・Lie環で現れるresultant零を厳密検算する。
滑らか群oidと幾何学的矢の全被覆は上の構造証明による。

\begin{theorem}\label{thm:h3-order2-full-inertia}
全$(1,3)$許容表示の幾何学的慣性は次の通りである。
\begin{enumerate}
\item $a=c=0$の閉部分は正常軸の二次元部分であり、慣性は$\mu_2\times\mu_2$。
\item $d=b=c=0$の閉部分は一次元である。
 $\bar g=g-ef/2,\ \bar h=h-f^2/4$と書くと許容条件は
 $\bar h(\bar g^2+e^2\bar h)\ne0$。
 慣性は$\bar g\ne0$で$\mu_4$、$\bar g=0$で$\mu_8$である。
\item その他の全点の慣性は$\mu_2$。
\end{enumerate}
第二部分の表示stackと粗空間はそれぞれ
\[
 [\operatorname{Spec}k[\kappa,(\kappa^2+1)^{-1}]/\mu_8],
 \qquad\mathbb A^1_j\setminus\{-1\},\quad
 j=\frac{\bar g^2}{e^2\bar h}=\kappa^2,
\]
ただし$\rho\in\mu_8$は$\kappa\mapsto\rho^{-4}\kappa$で作用する。
$\mu_8$点は一つの幾何学的同型類であり、
$(a,b,c,d,e,f,g,h)=(1,0,0,0,1,0,0,1)$を代表に取れる。
二つの慣性増大閉部分は互いに交わらない。
\end{theorem}
\begin{proof}
定理\ref{thm:h3-order2-full-arrows}でsourceとtargetを同一とし、
$\beta=1$で割る。$e\ne0$より$\gamma\delta=1$である。
$ad\ne0$なら最初と第四の式から$\alpha^2=\delta=\gamma=1$、
さらに$r=s=0$となるので慣性は$p$の符号だけである。

$a=0,d\ne0$では$v$の$u$平行移動で$b=0$にできる。
その後の自己同型では$r=s=0,\alpha^2=1$である。
$c\ne0$なら$\gamma^3=1$。
もし$\gamma\ne1$なら平方関係の$X,Y$係数から$g=h=0$となるが、
このとき$v^2=0$で整域性に反する。
$c=0$では正常軸部分となる。ここで$g=0$は許容されない：
判別四次は
\[
 (xz-dy^2-fx^2)(4ehx^2-ez^2)
\]
で、$[0:1:0]$は第一因子が非零の特異点だからrank2 fibre上の悪い判別点である。
従って$g\ne0$、自己同型式は$\delta^2=\gamma^2=1$となり、
$\alpha^2=1$と独立に二つの符号が残る。

$d=0,a\ne0$では$v$の$w$平行移動で$f=0$にする。
係数は$g\mapsto\bar g,h\mapsto\bar h$と変わる。
自己同型では$r=s=0,\delta=\alpha^2,\gamma=\alpha^{-2}$。
$b\ne0$なら$\delta=1$。
$b=0,c\ne0$なら$\delta^3=1$であり、$\delta\ne1$は再び$g=h=0$を強制する。
従って残る増大部分は$d=b=c=f=0$だけである。
ここではAS4条件と四次式が
\[
 \Delta_{\mathrm{AS}}=a^4(g^2+e^2h)\ne0,\qquad
 F_4=xz(4ehx^2+4gxz-ez^2)-e^2a^2y^4
\]
となる。$h\ne0$なら二項四次$xz(4ehx^2+4gxz-ez^2)$の四根は相異なる。
$\partial_yF_4=0$は$y=0$を強制するから$F_4$は滑らかである。
$h=0$なら$[1:0:0]$は退化した接錐を持つ特異点となり非許容である。
従って必要十分条件は定理の式である。
自己同型の残る条件は$g(\alpha^4-1)=0,h(\alpha^8-1)=0$で、
慣性は記載した巡回群となる。

この曲線上ではfppf局所的に$a=e=h=1$へ正規化できる。
実際、正規形への矢は$\gamma=e/\delta,\delta^4=he^2,\alpha^2=\delta/a$で与える。
残る$g'=g/\delta^2=\kappa$は$\kappa^2\ne-1$を満たす。
二つの正規形間の全矢は$\alpha^8=1$と
$\kappa'=\alpha^{-4}\kappa$だから、上の商stackと粗座標を得る。
正常軸部分の次元2は既存の全正常軸分類から、曲線の次元1はこの表示から従う。
両閉部分が交われば$a=d=0$となり許容係数条件に反する。
\end{proof}

\texttt{verify\_cubic\_order2\_full\_inertia.py}は全自己同型式、根の位数制約、
曲線上のresultant・判別式、除外する退化接錐とrank2特異点を厳密検算する。
これで$(1,3)$部分の全幾何学的慣性を決定したが、他スペクトルの被覆、
全H3の成分・交叉・族の内在的認識は依然として残る。
